% Copyright (c) 2026 Eugene Joseph M. Ragasa.
% SPDX-License-Identifier: Apache-2.0
% Author: Eugene Joseph M. Ragasa
% Formal proof completed: 2026-10-19
\documentclass[11pt]{article}
\usepackage[T1]{fontenc}
\usepackage{amsmath,amssymb,booktabs,tabularx}
\usepackage[margin=1in]{geometry}
\title{Proof Dossier: Exact Finite-Character Kernel}
\author{Eugene Joseph M. Ragasa}
\date{Formal proof completed: 2026-10-19}
\begin{document}\maketitle
\section{Identity and status}
\begin{tabularx}{\textwidth}{@{}lX@{}}\toprule
Statement ID & \texttt{BG-FOURIER-THM-009} \\
Status & proved; no admissions \\
Lean declaration & \texttt{ProjectKoios.BandGauge.Fourier.}\newline\texttt{finiteCharacterKernel} \\
Owning module & \texttt{ProjectKoios/BandGauge/Fourier/FiniteAbelianTransform.lean} \\
Axioms & \texttt{propext}, \texttt{Classical.choice}, \texttt{Quot.sound} \\
\bottomrule\end{tabularx}
\section{Statement}
For a finite abelian group $G$ and momentum labels $k,x\in G$,
\[
 \frac{1}{|G|}\sum_{\chi\in\widehat G}\chi(k)\overline{\chi(x)}
 =\begin{cases}1,&k=x,\\0,&k\ne x.\end{cases}
\]
\section{Argument}
Each summand is $\chi(k-x)$. Finite-character orthogonality evaluates its normalized sum as one exactly when $k-x=0$, which is equivalent to $k=x$.
\section{Boundary and provenance}
This standard finite harmonic-analysis identity is supporting infrastructure. It does not identify $G$ with a publication mesh or claim novelty. Lean checked the exact statement; the dossier and proof were AI-assisted.
\end{document}
